\section{Research Gaps and Directions}
\label{sec:gaps}

The final gap analysis will be derived from the validated evidence matrix rather than from the taxonomy alone. Candidate gaps are listed here only as working hypotheses and will be retained, revised, or rejected according to the final corpus.

\subsection{Quality Attribution Across Stages}
% Working hypothesis: rich component metrics but weak causal/cross-stage attribution.

\subsection{Reliability of Automated Evaluators}
% Working hypothesis: widespread judge use exceeds evidence on judge reliability.

\subsection{Trace and Process Visibility}
% Working hypothesis: final-response metrics miss important agent/process failures.

\subsection{Operational Quality as Part of Overall Quality}
% Working hypothesis: latency/cost/availability are often evaluated separately from
% semantic and task-level quality.

\subsection{Agreement Between Evaluation Frameworks}
% Working hypothesis: different frameworks/evaluators can disagree substantially on
% the same objects; investigate construct validity and calibration.

\subsection{From Measurement Collections to Integrated Quality Models}
% TODO: position the eventual PQI contribution only after evidence supports the gap.
% Do not freeze weights or a final mathematical formulation before synthesis.
